\chapter{TandemApp: the dashboard and proxy on their own}\label{ch:install-app}

The dashboard\index{dashboard} also runs without the engine, as TandemApp\index{TandemApp} (\code{tandem-app}), in front of any OpenAI-compatible engine or cloud API. Clients send their requests to one address; the proxy\index{proxy} routes each one by its model id to the backend that serves it, copies the answer back byte for byte, and records tokens, time to first token, decode speed, cached tokens and cost in one ledger. It needs Python 3.11 or newer and one package, \code{cryptography}. No GPU, no torch.

\section{Installing from a checkout}

\begin{shell}
git clone https://github.com/0xBakeer/TandemLLM.git && cd TandemLLM
pip install '.[host]'      # the host extra adds psutil and NVML for the System view
tandem-app                 # http://127.0.0.1:8080/dashboard/ opens on Setup
\end{shell}

The distribution is called \code{tandemllm}. Its \code{host} extra adds \code{psutil} and \code{nvidia-ml-py}, which the System view uses to read the machine; without them the collector falls back to \code{/proc} and \code{nvidia-smi}. The package is not on PyPI; install it from the checkout, or build the wheel with \code{python -m build --wheel}. \code{python -m proxy} in a checkout does the same as \code{tandem-app} without installing anything but \code{cryptography}.

The page opens on Setup\index{Setup screen} while no backend exists: add one, test the connection, save. The panels that only an engine fills (caches, speculation, the engine side of the conductor) say that this server has no engine.

\section{The tandem-app command}\index{tandem-app}

\begin{table}[htbp]
\centering\small
\caption{The subcommands of \texttt{tandem-app}.}\label{tab:tandem-app}
\begin{tabularx}{\linewidth}{@{}L{.3\linewidth}Y@{}}
\toprule
\hd{command} & \hd{what it does} \\
\midrule
\code{tandem-app}, \code{tandem-app serve} & the server: the proxy, the dashboard and the ledger; every flag of Table~\ref{tab:ref-proxy-flags} \\
\code{tandem-app check} & reads the environment and a copy of the config database through the schema and prints what would run: backends, key hints, locked fields, routing rules, the listen address. Writes nothing; exits 1 on an error that names the variable \\
\code{tandem-app secrets init} & makes the master key file beside the config database (mode 0600) when no master key answers; \code{secrets status} says which source answers \\
\code{tandem-app backends list} & one line per backend, from a copy of the database \\
\code{tandem-app hostcollect} & the remote host collector (Section~\ref{sec:hostcollect}) \\
\code{tandem-app launch vllm|sglang} & starts vLLM or SGLang with TandemLLM's plugin and registers it as a backend (Section~\ref{sec:launch}) \\
\code{tandem-app version} & the installed version \\
\bottomrule
\end{tabularx}
\end{table}

\section{A first backend}

A backend can be added in the GUI or declared in an env file. The file form suits a headless machine or a container. A local vLLM:

\begin{shell}
cat > app.env <<'EOF'
TANDEM_BACKENDS=vllm1
TANDEM_BACKEND_VLLM1_KIND=vllm
TANDEM_BACKEND_VLLM1_URL=http://127.0.0.1:8100/v1
EOF
tandem-app check --env-file app.env        # prints what would run, writes nothing
tandem-app serve --env-file app.env
\end{shell}

A cloud API, with the key in the environment, where it stays in memory and is never written to the database:

\begin{shell}
printf 'TANDEM_BACKENDS=deepseek\nTANDEM_BACKEND_DEEPSEEK_KIND=deepseek\nTANDEM_BACKEND_DEEPSEEK_KEY=%s\n' \
    "$DEEPSEEK_API_KEY" > app.env && chmod 600 app.env
tandem-app serve --env-file app.env &
curl -s localhost:8080/v1/chat/completions -H 'Content-Type: application/json' \
    -d '{"model": "deepseek-flash", "messages": [{"role": "user", "content": "Say hi"}]}'
\end{shell}

Then open \code{http://127.0.0.1:8080/dashboard/#/usage}: the request is one row there, with its tokens, its time to first token and its cost at the model's price. Chapter~\ref{ch:other-engines} has a section for each kind of backend.

\section{In docker}\index{docker}

\code{ops/Dockerfile.app} builds an image on \code{python:3.12-slim} from the same wheel, with the host extra and the built page inside, and no compiler or torch. It is configured by an env file; \code{ops/app.env.example} is one to start from.

\begin{shell}
docker build -f ops/Dockerfile.app -t tandemllm-app .
cp ops/app.env.example app.env            # fill in; no quotes around values
docker run -d --name tandem-app -p 127.0.0.1:8080:8080 -v tandem-app:/data \
    --env-file app.env tandemllm-app
\end{shell}

The config database, the usage ledger and the key file live in the \code{/data} volume. Inside the container the server binds \code{0.0.0.0}, so what the outside reaches is what \code{-p} publishes: \code{127.0.0.1:8080} for this machine only. Keep it that way: no admin password is set by default, and this release has no sign-in screen (Section~\ref{sec:access}). The image runs as an unprivileged user and answers its health check on \code{/health}. The release's tests build and install the wheel; they do not build the image. A volume from 0.15.0 or before keeps its old name; \code{docs/proxy.md} says how to mount or copy it.

Inside a container the host collector sees the container, not the machine. Run it with \code{--pid=host} (and \code{--gpus all} for NVML) for the machine's view, or run \code{tandem-app hostcollect}\index{tandem-app hostcollect@\texttt{tandem-app hostcollect}} on the machine itself.

\section{Where the proxy keeps its files}

\begin{table}[htbp]
\centering\small
\caption{The standalone proxy's files.}\label{tab:proxy-files}
\begin{tabularx}{\linewidth}{@{}L{.18\linewidth}YL{.36\linewidth}@{}}
\toprule
\hd{file} & \hd{what} & \hd{default} \\
\midrule
config database & backends, encrypted keys, models, prices, routing rules, users, the fleet's nodes and saved window layouts, an audit table & \code{~/.tandemllm/config.sqlite3} (\code{--config-db}, \code{TANDEM_CONFIG_DB}) \\
usage ledger & one row per request: tokens, times, cost; no text & \code{~/.tandemllm/ledger.sqlite3} (\code{--usage-ledger}; \code{off} for none) \\
master key file & the key the backends' keys are encrypted with, when no other source answers & \code{~/.tandemllm/secrets.key}, mode 0600 \\
conductor documents & the harness's profiles, router, packages and journal & \code{~/.tandemllm/conductor} (\code{--conductor-config}; \code{off} for none) \\
\bottomrule
\end{tabularx}
\end{table}

\section{The proxy inside the engine's server}

The engine's own server can relay too. Start it with \code{--backends <config database>} (\code{QSE_BACKENDS}\index{qsebackends@\texttt{QSE\_BACKENDS}} in \code{serve.env}; empty means off). It then serves its own model as before and relays every request whose model belongs to a backend. A relayed request runs on its handler's own thread and never takes the engine's lock or a place in its queue, so a slow cloud answer cannot hold up the engine, and the engine's queue cannot hold up a cloud answer. With \code{--backends}\index{backends@\texttt{-\kern0pt-backends}} off the engine loads none of the proxy's code and serves exactly as without it.

\section{The fleet}\label{sec:fleet}\index{fleet}

TandemApp keeps the registry of the Sparks it manages and answers the fleet API, \code{/v1/fleet/*}, which the dashboard's shell, scope bar and windows read. The contract is \code{docs/contract/fleet-v1}. The browser talks to \code{tandem-app} only: every figure of a node is read on the server.

On one box the fleet has one node: the machine \code{tandem-app} runs on, when it has an NVIDIA GPU or an engine of its own. The first start registers it with its hostname as the alias, the first identity colour and an id from the machine (the DMI serial, else the machine id), never from its address; a new address keeps the node and is noted in the log. A machine in front of cloud APIs only, a Mac for instance, is the controller's host and no node. A node on another machine is listed as off until its node agent and pairing exist.

\begin{table}[htbp]
\centering\small
\caption{The fleet API on a fleet of one.}\label{tab:fleet-one}
\begin{tabularx}{\linewidth}{@{}L{.26\linewidth}Y@{}}
\toprule
\hd{route} & \hd{what it answers} \\
\midrule
\code{nodes}, \code{nodes/<id>} & the node: status, version and inventory (OS, kernel, driver, CUDA, ConnectX ports, disk, interfaces); \code{PATCH} sets the alias, the colour (1 to 8, one node each), the icon and the tags \\
\code{deployments} & each model an engine on the node serves itself, as a Single deployment, its status from the backend's health check \\
\code{links}, \code{topology} & the LAN of the default route, and the internet when a cloud backend exists \\
\code{metrics} & decode tok/s, requests and TTFT from each engine's \code{/metrics}; GPU, power, temperature, clocks, memory, CPU, disk, LAN and fabric from the host collector. Sampled every 2\,s into memory: 10 minutes at 2\,s, 7 days at 30\,s. History starts with the process; a point without a sample is null \\
\code{memory} & the engine's split (weights, KV, caches) beside the rest of the machine's use and the page cache, GiB \\
\code{logs} & \code{tandem-app}'s own log and its TandemLLM engines', merged by time \\
\code{layouts/<page>} & a page's window layout per device class, in the config database; the last write wins \\
\code{events}, \code{alerts}, \code{rules} & empty lists until the fleet's event log and alert rules exist \\
\bottomrule
\end{tabularx}
\end{table}

Reads follow the dashboard rule and writes the admin rule (Section~\ref{sec:access}), as for \code{/v1/config}. The engine's own server, which serves the dashboard on a box without \code{tandem-app}, has no fleet API; the dashboard then shows the box alone.

\section{Access}

The standalone proxy follows the engine's access rule (Section~\ref{sec:access}) with its own variable, \code{TANDEM_ADMIN_TOKEN}, and adds user keys for the OpenAI routes (Section~\ref{sec:users}). It binds \code{127.0.0.1} by default.
